\chapter{Contexto}
\label{chap:contexto}

En este capítulo se establece el contexto general en el que se enmarca el presente Trabajo de Fin de Grado (TFG). En primer lugar, se exponen los antecedentes, definiendo la problemática actual. Posteriormente, se detalla la motivación fundamental del trabajo y se establecen los objetivos, tanto generales como específicos, que guiarán la investigación y el desarrollo de este proyecto.

\section{Antecedentes}
\label{sec:antecedentes}

TODO

\section{Motivación y objetivos}
\label{sec:motivacion_objetivos}

\subsection{Motivación}
Aunque las herramientas actuales de gestión han integrado con éxito visualizaciones clásicas, estas en su mayoría perpetúan paradigmas analógicos. Esta herencia visual desaprovecha las capacidades de las interfaces modernas para sintetizar la complejidad, lo que genera dos barreras principales en la industria actual:

\begin{itemize}
    \item \textbf{Alta carga cognitiva y fatiga visual}: Para comprender el estado real de un proyecto, los gestores y desarrolladores deben procesar manualmente una gran cantidad de texto y datos abstractos dispersos. La falta de representaciones visuales intuitivas dificulta la detección rápida de bloqueos, dependencias o riesgos, obligando al usuario a realizar un esfuerzo mental considerable para construir un modelo mental del proyecto.
    \item \textbf{Burocracia frente a trabajo de valor}: Gran parte del tiempo de gestión se invierte en tareas mecánicas y repetitivas: solicitar actualizaciones de estado, requerir la cumplimentación de campos y sintetizar información fragmentada. Esta microgestión interrumpe los estados de concentración de los perfiles técnicos y desvía a los responsables de la toma de decisiones estratégicas.
\end{itemize}

\subsection{Objetivos}

\subsubsection{Objetivo General}
Investigar, diseñar y desarrollar un nuevo paradigma de plataforma de gestión de proyectos que reemplace las vistas tabulares convencionales por metáforas visuales inmersivas, capaces de representar la complejidad estructural del trabajo. Asimismo, se pretende concebir e integrar un sistema inteligente capaz de actuar de forma autónoma para asistir en la coordinación asíncrona y el seguimiento de los equipos.

\subsubsection{Objetivos Específicos}
Para alcanzar el objetivo general, se plantean los siguientes objetivos específicos:

\begin{itemize}
    \item \textbf{A. Investigación y análisis}: Analizar las limitaciones de usabilidad y experiencia de usuario (UX) en las herramientas de gestión convencionales y estudiar su impacto en la carga cognitiva. Investigar patrones de diseño de interfaces en sectores ajenos a la gestión empresarial que permitan representar jerarquías y dependencias de forma intuitiva.
    
    \item \textbf{B. Diseño conceptual y visual}: Definir una propuesta de interfaz gráfica y experiencia de usuario que permita visualizar la topología de un proyecto como un sistema interconectado. El objetivo es priorizar la claridad visual frente a la lectura extensiva de texto, fomentando una interacción más natural y adaptada a entornos profesionales.
    
    \item \textbf{C. Arquitectura de agentes inteligentes}: Diseñar un sistema basado en Inteligencia Artificial capaz de actuar como intermediario ágil entre el usuario y la plataforma. El propósito es establecer flujos de trabajo que permitan delegar tareas repetitivas de coordinación, comunicación asíncrona y seguimiento de estados.
    
    \item \textbf{D. Ingeniería y arquitectura de software}: Diseñar e implementar una base técnica robusta que garantice la escalabilidad, el alto rendimiento y la mantenibilidad del código. Se persigue estructurar el sistema aplicando principios de diseño de software avanzados que aseguren un alto grado de desacoplamiento entre los dominios de la aplicación, facilitando su evolución y mantenimiento futuros.
    
    \item \textbf{E. Prototipado y validación técnica}: Desarrollar un Producto Mínimo Viable (MVP) que integre los nuevos paradigmas de visualización y el sistema multi-agente. Validar la viabilidad de la arquitectura propuesta y evaluar si la solución aporta un valor funcional tangible frente a los enfoques de gestión de proyectos tradicionales.
\end{itemize}